
This section enumerates the risks associated with the protocol, the controls in place, and residual risks that users must accept. Each risk category follows a consistent schema: threat, impact, preconditions, controls, and residual risk.

\paragraph{Threat model assumptions.}
\begin{itemize}
\item The Starknet network operates correctly (liveness, validity proofs, L1 finality).
\item The Pragma oracle infrastructure is not fundamentally compromised.
\item Cryptographic primitives (hashing, signatures) remain secure.
\item Users can monitor on-chain events and act within timelock windows.
\end{itemize}

\section{Smart Contract Risk}

\subsection{Code Vulnerabilities}

\begin{center}
\begin{tabularx}{\textwidth}{l >{\raggedright\arraybackslash}X l l}
\toprule
\textbf{Threat} & \textbf{Controls} & \textbf{Status} & \textbf{Residual} \\
\midrule
Reentrancy & Checks-effects-interactions; explicit reentrancy guards; minimize external calls in state-mutating paths & Implemented & Low \\
Arithmetic overflow & Widened-precision math; explicit overflow checks; fuzz testing over boundary values & Implemented & Low \\
Access control bypass & Role-based access control (RBAC); explicit permission checks; separation of duties & Implemented & Low \\
Logic errors & Unit tests; integration tests; invariant fuzzing; external audit & Partial & Medium \\
Front-running / MEV & User-defined slippage bounds; min\_out / max\_in parameters & Implemented & Medium \\
\bottomrule
\end{tabularx}
\end{center}

\paragraph{Note on Cairo.} Starknet's execution model differs from EVM but does not eliminate reentrancy-like patterns through external calls and callbacks. Explicit guards and patterns are required.

\subsection{Dependency Risk}

\begin{center}
\begin{tabular}{lll}
\toprule
\textbf{Dependency} & \textbf{Controls} & \textbf{Residual Risk} \\
\midrule
OpenZeppelin Cairo & Pinned versions; monitor security advisories & Upstream vulnerability \\
cairo-fp (fixed-point math) & Extensive precision testing; boundary validation & Math library bugs \\
Starknet protocol & Accept as infrastructure dependency & Protocol-level bugs \\
\bottomrule
\end{tabular}
\end{center}

\section{Upgrade Risk}

All core contracts are upgradeable using OpenZeppelin's Upgradeable component (class hash replacement). This section specifies the upgrade authority, controls, and user protections.

\subsection{Upgrade Impact by Contract}

\begin{center}
\begin{tabular}{lll}
\toprule
\textbf{Contract} & \textbf{Upgrade Can Modify} & \textbf{Impact} \\
\midrule
SY & Deposit/redeem logic; share accounting & Critical \\
PT & Redemption logic; mint authority & Critical \\
YT & Yield distribution; interest accounting & Critical \\
Market & AMM curve; swap logic; fee calculation & Critical \\
Router & Routing logic; could add hidden fees & High \\
Factory / MarketFactory & Implementation templates for new deployments & High \\
\bottomrule
\end{tabular}
\end{center}

\subsection{Upgrade Policy}

\begin{center}
\begin{tabularx}{\textwidth}{l l >{\raggedright\arraybackslash}X l}
\toprule
\textbf{Control} & \textbf{Status} & \textbf{Parameter} & \textbf{Residual} \\
\midrule
Multisig ownership & Planned & Target: 3-of-5 signers & Collusion risk \\
Timelock delay & Planned & Target: 48 hours minimum & Emergency bypass \\
Upgrade events & Implemented & \texttt{Upgraded(class\_hash)} emitted & Requires monitoring \\
Code verification & Required & All class hashes verified on-chain & None \\
Pre-deployment audit & Policy & Audited before mainnet deployment & Audit scope limits \\
\bottomrule
\end{tabularx}
\end{center}

\paragraph{Emergency upgrade policy.}
\begin{itemize}
\item Emergency upgrades (bypassing timelock) are permitted only for active exploits.
\item Emergency actions require documented justification published within 24 hours.
\item Post-emergency review and standard re-deployment within 7 days.
\end{itemize}

\subsection{User Protections During Upgrades}

\begin{center}
\begin{tabularx}{\textwidth}{l >{\raggedright\arraybackslash}X}
\toprule
\textbf{Guarantee} & \textbf{Description} \\
\midrule
Exit window & Timelock provides users N hours to exit before upgrade execution. \\
Always-available operations & Even during pause: \texttt{redeem\_pt\_after\_expiry}, \texttt{burn} (LP withdrawal), \texttt{claim\_interest}. \\
Pausable operations & Swaps and new mints may be paused; existing positions remain accessible. \\
Upgrade transparency & Upgrade proposals include class hash diff and audit link (policy). \\
\bottomrule
\end{tabularx}
\end{center}

\section{Oracle Risk}

The protocol depends on Pragma oracle for yield index data used in YT interest accounting.

\subsection{Oracle Parameters}

\begin{center}
\begin{tabular}{ll}
\toprule
\textbf{Parameter} & \textbf{Value} \\
\midrule
TWAP window & 1 hour (configurable per SY) \\
Staleness threshold & 86400 seconds (24 hours) \\
Data source & Pragma aggregated feed \\
\bottomrule
\end{tabular}
\end{center}

\subsection{Oracle Failure Behavior}

\begin{center}
\begin{tabularx}{\textwidth}{l l >{\raggedright\arraybackslash}X}
\toprule
\textbf{Condition} & \textbf{Protocol Response} & \textbf{Rationale} \\
\midrule
Stale data (age $>$ threshold) & Use watermark index; no new yield distributed & Protects against stale/manipulated data \\
Oracle unavailable & Operations using index continue with last value & Prevents liveness failure \\
Rate below watermark & Watermark unchanged; yield accounting frozen & Prevents negative yield distribution \\
\bottomrule
\end{tabularx}
\end{center}

\paragraph{Operations affected by oracle state.}
\begin{itemize}
\item \texttt{mint\_py}, \texttt{redeem\_py}, \texttt{claim\_interest}: Use current or watermark index.
\item \texttt{redeem\_pt\_after\_expiry}: Does not depend on oracle (1:1 redemption).
\item AMM swaps: Do not use oracle; pricing is internal to pool state.
\end{itemize}

\subsection{Exchange Rate Manipulation}

\paragraph{Threat.} Attacker manipulates oracle to extract value from the protocol.

\paragraph{Attack surfaces.} The oracle index affects only the YT interest accounting system:

\begin{center}
\begin{tabularx}{\textwidth}{l >{\raggedright\arraybackslash}X l}
\toprule
\textbf{Subsystem} & \textbf{Oracle Dependency} & \textbf{Extraction Vector} \\
\midrule
Interest accounting & $I_{global}$ updated from oracle & Inflated index $\Rightarrow$ excess \texttt{claim\_interest} payout \\
SY deposit/redeem & 1:1 ratio (no oracle) & None \\
AMM pricing & Internal pool state only & None (oracle not used) \\
\bottomrule
\end{tabularx}
\end{center}

\paragraph{Attack scenario (index inflation).}
\begin{enumerate}
\item Attacker manipulates oracle to report inflated exchange rate.
\item Attacker calls \texttt{claim\_interest}; inflated $I_{global}$ causes excess SY payout.
\item Oracle returns to normal; attacker has extracted SY from YT contract reserves.
\end{enumerate}

\paragraph{Controls.}
\begin{itemize}
\item TWAP (5-minute window) increases manipulation cost and duration required.
\item Watermark prevents exploitation of rate \emph{drops} (only increases can inflate index).
\item Staleness check rejects data older than 24 hours.
\end{itemize}

\paragraph{Residual risk.} Sustained oracle manipulation over the TWAP window can still inflate the index. Mitigation requires oracle-level security (Pragma's aggregation and data source diversity).

\section{Market Risk}

Market risks affect execution and returns but do not compromise protocol solvency. These are inherent to participation in yield markets.

\subsection{Liquidity Risk}

\begin{center}
\begin{tabularx}{\textwidth}{l >{\raggedright\arraybackslash}X l}
\toprule
\textbf{Scenario} & \textbf{Impact} & \textbf{User Mitigation} \\
\midrule
Thin liquidity & High slippage; inability to exit at fair price & Smaller orders; check price impact \\
LP concentration & Single LP can move market significantly & Monitor LP distribution \\
Extreme proportion & Trades rejected at bounds & Wait for arbitrage; use redemption path \\
\bottomrule
\end{tabularx}
\end{center}

\paragraph{Worst case.} Inability to exit position at fair price before expiry. Protocol solvency is unaffected; this is execution risk only.

\subsection{Impermanent Loss Risk}

IL risk for LPs depends on rate volatility and time to expiry:

\begin{itemize}
\item \textbf{High volatility + long maturity}: Severe IL exposure; rate movements cause significant rebalancing.
\item \textbf{Low volatility or short maturity}: Minimal IL; PT price convergence dominates.
\item \textbf{Near expiry}: IL minimizes as $P_{PT} \to 1$ regardless of rate movements.
\end{itemize}

\paragraph{Key insight.} IL risk decreases as expiry approaches because PT price becomes insensitive to rate changes ($r_{scalar} \to \infty$).

\section{Underlying Asset Risk}

Risks from the yield-bearing assets wrapped by SY tokens.

\begin{center}
\begin{tabularx}{\textwidth}{l >{\raggedright\arraybackslash}X >{\raggedright\arraybackslash}X}
\toprule
\textbf{Risk} & \textbf{Protocol Behavior} & \textbf{User Impact} \\
\midrule
Depeg (value loss) & Watermark prevents negative yield accounting; redemption continues at 1:1 in SY & Principal impairment; users redeem into devalued asset \\
Yield cessation & YT stops accruing; PT converges toward 1 SY & YT holders receive no further yield; PT may reprice \\
Underlying protocol failure & SY may become worthless & Total loss possible \\
Regulatory freeze & Protocol cannot intervene & User bears external legal risk \\
\bottomrule
\end{tabularx}
\end{center}

\paragraph{Key distinction.} The watermark protects yield accounting integrity (no negative yield distribution). It does \emph{not} protect against principal impairment of the underlying asset. If the underlying loses value, all token holders (SY, PT, YT, LP) bear proportional losses.

\section{Governance and Admin Risk}

\subsection{Admin Key Compromise}

\paragraph{Threat.} Attacker gains control of admin keys and executes malicious upgrades.

\paragraph{Impact.} Complete protocol control; ability to drain funds via malicious upgrade.

\begin{center}
\begin{tabularx}{\textwidth}{l l >{\raggedright\arraybackslash}X}
\toprule
\textbf{Control} & \textbf{Status} & \textbf{Protection} \\
\midrule
Hardware wallet storage & Required & Reduces remote key theft \\
Multisig (3-of-5) & Planned & Requires multiple compromises or collusion \\
Timelock (48h) & Planned & Users can exit before malicious upgrade executes \\
Upgrade monitoring & Implemented & On-chain events enable automated alerts \\
\bottomrule
\end{tabularx}
\end{center}

\subsection{Malicious Upgrade Detection}

\begin{center}
\begin{tabularx}{\textwidth}{l >{\raggedright\arraybackslash}X >{\raggedright\arraybackslash}X}
\toprule
\textbf{Attack} & \textbf{Detection Method} & \textbf{User Defense} \\
\midrule
Backdoor insertion & Code diff analysis; class hash verification & Exit during timelock window \\
Hidden fee extraction & Event monitoring; balance tracking & Review upgrade proposals \\
Logic manipulation & Automated contract comparison tools & Subscribe to upgrade alerts \\
\bottomrule
\end{tabularx}
\end{center}

\subsection{Role Escalation}

\begin{center}
\begin{tabular}{ll}
\toprule
\textbf{Risk} & \textbf{Control} \\
\midrule
Admin grants attacker roles & Multisig requirement; role changes emit events \\
Role accumulation & Separation of duties; distinct admin/pauser/operator roles \\
Pauser griefing & Pause is temporary; unpause requires same authority \\
\bottomrule
\end{tabular}
\end{center}

\section{Starknet Infrastructure Risk}

External infrastructure risks outside protocol control.

\begin{center}
\begin{tabularx}{\textwidth}{l >{\raggedright\arraybackslash}X >{\raggedright\arraybackslash}X}
\toprule
\textbf{Risk} & \textbf{Impact} & \textbf{User Response} \\
\midrule
Sequencer downtime & Transactions not processed & Wait for recovery; funds remain in L2 state \\
Sequencer censorship & Specific transactions blocked & Use alternative RPC; await escape mechanisms \\
Proof generation delay & Finality delayed on L1 & Internal accounting unaffected; settlement certainty reduced \\
L1 congestion & High finalization costs & Accept as infrastructure cost \\
Protocol upgrade (Starknet) & Breaking changes possible & Monitor Starknet roadmap; protocol will adapt \\
\bottomrule
\end{tabularx}
\end{center}

\paragraph{Note on censorship.} L1 force-inclusion mechanisms for Starknet are evolving. Until mature escape hatches exist, sequencer censorship remains an infrastructure-level risk that users accept.

\section{User Operational Risk}

User errors and operational security failures.

\begin{center}
\begin{tabularx}{\textwidth}{l >{\raggedright\arraybackslash}X >{\raggedright\arraybackslash}X}
\toprule
\textbf{Risk} & \textbf{Description} & \textbf{Mitigation} \\
\midrule
Wrong address & Funds sent to incorrect recipient & Address validation; UI confirmations; verify full address \\
Address poisoning & Attacker creates similar-looking address & Use address book; verify complete address, not prefix \\
Wrong network/contract & Interaction with fake or test contract & Verify class hash; use official registry \\
Missed expiry & PT held past maturity & Redemption still works; no penalty \\
Slippage misconfiguration & Worse execution than expected & UI defaults; warnings on extreme values \\
Unlimited approvals & Excessive token permissions & Default to exact approvals; revoke after use \\
Approval phishing & Malicious approval requests & Wallet warnings; review approval targets \\
MEV sandwich & Execution between attacker trades & Tight slippage bounds; private mempools if available \\
Private key loss & Permanent fund loss & Hardware wallets; secure backup procedures \\
\bottomrule
\end{tabularx}
\end{center}

\section{Threat Model Summary}

\subsection{Risk Priority Matrix}

Risks are prioritized separately for systemic (protocol-wide) and individual (user-level) impact.

\begin{center}
\begin{tabularx}{\textwidth}{l l l l l}
\toprule
\textbf{Risk} & \textbf{Likelihood} & \textbf{Systemic Impact} & \textbf{User Impact} & \textbf{Priority} \\
\midrule
Admin key compromise & Low & Critical & Critical & Critical \\
Smart contract bug & Low & Critical & Critical & Critical \\
Oracle manipulation & Medium & High & High & High \\
Underlying depeg & Low & High & High & High \\
Liquidity crisis & Medium & Low & Medium & Medium \\
Sequencer downtime & Low & Low & Low & Low \\
User operational error & High & None & High & Medium (user) \\
\bottomrule
\end{tabularx}
\end{center}

\paragraph{Interpretation.}
\begin{itemize}
\item \textbf{Critical priority}: Requires maximum preventive investment and monitoring.
\item \textbf{High priority}: Active controls required; residual risk accepted.
\item \textbf{Medium priority}: Controls in place; ongoing monitoring.
\item \textbf{Low priority}: Accepted as infrastructure or user-level risk.
\end{itemize}

\section{Exclusions: What the Protocol Does NOT Protect Against}

The following risks are explicitly outside the protocol's protection scope:

\begin{enumerate}
\item \textbf{Underlying asset impairment.} Principal loss from underlying asset failure is borne by users.

\item \textbf{Malicious or compromised governance.} Upgrade authority can modify contract logic; mitigated only by timelock and monitoring where implemented.

\item \textbf{Starknet network failures.} Protocol correctness depends on Starknet liveness and validity proofs.

\item \textbf{Oracle compromise.} Fundamental Pragma compromise could cause incorrect yield accounting; mitigated by TWAP and watermark but not eliminated.

\item \textbf{Regulatory actions.} Legal restrictions on assets or users cannot be prevented by the protocol.

\item \textbf{User operational errors.} Incorrect addresses, lost keys, and configuration mistakes are user responsibility.

\item \textbf{MEV and front-running.} Users can mitigate with slippage settings but cannot eliminate; this is inherent to public blockchains.
\end{enumerate}
